\section{Vấn đề cần giải quyết}
\label{sec:problem}
\subsection{Bối cảnh và động cơ}
Một agent sử dụng mô hình ngôn ngữ lớn (LLM) thường lặp lại ba bước: đọc yêu cầu và dữ liệu, đề xuất lời gọi công cụ, rồi dùng kết quả công cụ để quyết định bước tiếp theo. Khác với hệ thống chỉ sinh văn bản, lời gọi có thể gửi thư, thêm lịch hoặc sửa tệp. Vì vậy, sai lệch trong suy luận có thể dẫn đến thay đổi thực tế ở môi trường.

Prompt injection gián tiếp xuất hiện khi chỉ dẫn của bên thứ ba nằm trong dữ liệu mà agent được phép đọc, thay vì trong yêu cầu trực tiếp của người dùng. Ví dụ, một tài liệu công việc có thể chứa nội dung yêu cầu agent thực hiện hành động ngoài nhiệm vụ. Ranh giới cần bảo vệ là quyền hành động: việc đọc được một chỉ dẫn không đồng nghĩa với việc được phép thực hiện chỉ dẫn đó \cite{greshake2023,debenedetti2024}.

Hệ thống còn phải đối mặt với thay đổi giao diện công cụ. Policy có thể kiểm tra một trường tham số ở phiên bản cũ nhưng không còn diễn giải đúng cấu trúc ở phiên bản mới. Một thay đổi vẫn hợp lệ về schema chưa chắc tương đương về quyền truy cập. Kiểm thử hồi quy cần quan sát cả tác động bị cấm lẫn khả năng hoàn thành tác vụ hợp lệ.

\subsection{Câu hỏi và mục tiêu}
Câu hỏi nghiên cứu là: một phòng thủ đã ngăn tác động có hại ở hợp đồng công cụ $v0$ có duy trì được sự bảo vệ khi chuyển sang $v1$, và việc chuyển đổi policy có khắc phục sai lệch mà vẫn giữ utility hay không? Trong phạm vi học phần, nhóm xây dựng thiết kế đo lường và hiện thực một phần giải pháp để kiểm chứng khả năng triển khai.

Ba đầu ra kỹ thuật của báo cáo là: định nghĩa so sánh theo cặp; xây dựng guard giới hạn quyền theo nhiệm vụ; và kiểm tra độc lập trạng thái sau thực thi. Đề tài cố định có tên tiếng Anh là \emph{Security Regression Testing of Indirect Prompt Injection Defenses under Versioned Agent Tool Contracts}. Báo cáo không khẳng định đã tạo một cơ chế phòng thủ mới hoặc đã chứng minh hiệu quả trên một quần thể agent.

\subsection{Phạm vi và mô hình đe dọa}
Bên tấn công có thể kiểm soát nội dung email, tài liệu hoặc kết quả công cụ mà agent đọc. Họ không được sửa policy tin cậy, dispatcher, mã công cụ hoặc oracle. Môi trường triển khai là sandbox với dữ liệu giả lập. Phạm vi không bao gồm đầu độc mô hình, chiếm quyền mã Python trong tiến trình tin cậy, hoặc bảo đảm luồng thông tin xuyên nhiều agent. Các giả định này giới hạn ý nghĩa của kết quả: guard tại dispatcher không bảo vệ một đường thực thi cố tình bỏ qua dispatcher.
